\label{ch:introduction}

This chapter introduces the purpose and overall direction 
of the thesis. It presents the motivation for the study 
and describes the main problem that is addressed. 
The research questions are then defined to guide the work. 
Finally, the chapter provides an overview 
of the structure of the thesis.

\section{Motivation}
\label{sec:motivation}

Machine learning has seen strong growth in recent years 
and is now widely applied in product development 
and production. It is used for tasks such as optimization, 
control, prediction, and decision support 
\parencite{wuest2016ml_manufacturing}. 

Advances in data availability and computing power 
have further increased the use of machine learning 
in industrial settings \parencite{rai2021ml_manufacturing_industry4}. 
As a result, machine learning plays an important role 
in improving efficiency, supporting decision-making, 
and enabling new capabilities in engineering and manufacturing.

At the same time, the field of machine learning 
is large and diverse. Many different algorithms, 
methods, and approaches are available, 
each with its own strengths and weaknesses 
\parencite{wuest2016ml_manufacturing}. 
There is no single method that performs best 
for all problems, and performance depends on 
the specific application and data 
\parencite{adekoya2022applications_ai_em,lee2020machine_learning_enterprises}. 

This creates a challenge for engineers and practitioners, 
who must select suitable methods for their problems. 
The large number of available options can act as a barrier 
and limit the effective use of machine learning in practice 
\parencite{adekoya2022applications_ai_em}. 
Selecting an appropriate method often requires balancing 
multiple factors such as accuracy, interpretability, 
and data requirements 
\parencite{plathottam2023review_ai_manufacturing,ali2017accurate_multi_criteria_decision_making_methodology_recommending_machine_learning_algorithm}. 

If the selection is not done carefully, the resulting model 
may perform poorly or fail to provide useful insights 
\parencite{ali2017accurate_multi_criteria_decision_making_methodology_recommending_machine_learning_algorithm}. 
In practice, simpler models are often preferred 
due to their interpretability, even if more complex models 
may offer higher accuracy \parencite{paleyes2022challenges_deploying_ml_survey_case_studies}. 
This highlights the importance of considering trade-offs 
in real-world applications.

Another challenge is the gap between research 
and practical implementation. Research results are 
distributed across many articles and domains, 
which makes it difficult to gain a complete overview 
\parencite{adekoya2022applications_ai_em}. 
In addition, most studies focus on developing models, 
while less attention is given to how they should be selected 
and applied in practice.

Companies also face challenges related to lack of knowledge 
and experience when adopting machine learning 
\parencite{sonntag2023pattern_ml_product_development}. 
In many cases, the selection of methods is based on 
trial and error, which is time-consuming and inefficient 
\parencite{sonntag2023pattern_ml_product_development}. 
There is also often pressure to adopt machine learning 
without a clear plan, which increases the risk 
of unsuccessful implementations 
\parencite{plathottam2023review_ai_manufacturing}. 

These challenges highlight the need for more structured 
approaches that can support method selection 
and make machine learning more accessible.

% \begin{itemize}
%     \item Increasing use of machine learning in product development and production
%     \item Difficulty in selecting suitable ML methods for specific problems
%     \item Gap between research literature and practical implementation
%     \item Need for structured decision support
%     \item Potential value of combining ontology and literature-based retrieval
%     \item Relevance for engineers and practitioners with limited ML expertise
% \end{itemize}

\section{Problem Description}
\label{sec:problem-description}

Despite the growing adoption of machine learning in manufacturing 
and engineering, practitioners often lack adequate support when 
selecting appropriate methods for their specific problems. 
Although a large number of use cases exist, there is currently 
no guidance or standard for developing machine learning solutions 
from ideation to deployment, and it remains difficult for 
non-experts to begin developing solutions for their specific 
problems \parencite{chen2023ml_manufacturing_industry4}. 
In practice, method selection is often based on trial and error, 
which is both time-consuming and dependent on prior experience 
that many practitioners do not have 
\parencite{sonntag2023pattern_ml_product_development}.

A central limitation is that no structured approach currently 
exists for making an informed selection of machine learning 
algorithms based on a well-defined problem description 
\parencite{sonntag2023pattern_ml_product_development}. 
Existing frameworks tend to use accuracy or prediction speed 
as the primary selection criteria, which are often too general 
and can lead to unsuitable methods being chosen 
\parencite{sala2018selecting_ml_algorithm}. 
Manual evaluation of candidate algorithms is highly time-consuming, 
and if the selection is not done carefully, 
the resulting model may fail to produce useful results 
\parencite{ali2017accurate_multi_criteria_decision_making_methodology_recommending_machine_learning_algorithm}. 

Existing research further focuses primarily on the development 
and evaluation of machine learning models, while less attention 
is given to what conditions lead to selecting one approach 
over another in specific industrial contexts 
\parencite{arena2024conceptual_framework_ml_algorithm_selection_predictive_maintenance}. 
The lack of interpretability in many approaches is also a key 
barrier to adoption, as practitioners need to understand 
and trust the recommendations they receive 
\parencite{presciuttini2024ml_iot_manufacturing_operations}.

At the same time, the relevant knowledge is distributed across 
a large and fragmented body of literature. 
The extensive number of works on machine learning algorithms 
and applications creates broad knowledge on the topic, but may 
also generate disorientation when selecting the right approach 
\parencite{sala2018selecting_ml_algorithm}. 
This knowledge is difficult to access and apply systematically 
during method selection, which limits the ability to make 
informed, evidence-based decisions 
\parencite{adekoya2022applications_ai_em}.

Addressing these challenges requires a structured approach 
that can translate problem descriptions into method 
recommendations and connect those recommendations to 
relevant research. 
This thesis investigates whether combining ontology-based 
knowledge representation with semantic retrieval from 
scientific literature can provide such support, offering 
interpretable recommendations grounded in prior research 
alongside guidance for early-stage implementation.

\textbf{Burde MLguide forklares i så stor detalj? Og bør det forklares her eller under research questions?}

% \begin{itemize}
%     \item Users struggle to identify relevant ML methods for their problem
%     \item Existing approaches are either too generic or too technical
%     \item Lack of tools linking:
%     \begin{itemize}
%         \item problem descriptions
%         \item ML methods
%         \item supporting research articles
%     \end{itemize}
%     \item Challenge of utilizing large volumes of literature effectively
%     \item Need for:
%     \begin{itemize}
%         \item interpretable recommendations
%         \item traceability to literature
%         \item support for early implementation steps
%     \end{itemize}
%     \item Context: system combining ontology + embeddings + article retrieval
% \end{itemize}

\section{Research Questions}
\label{sec:research-questions}

The challenges outlined above point to a need for structured 
decision support that can connect problem descriptions with 
suitable machine learning methods, use the research literature 
to justify those recommendations, and lower the barrier to 
early-stage implementation. This thesis investigates whether 
a system combining these capabilities can provide useful and 
justified recommendations in practice. The following hypothesis 
is proposed:


\begin{quote}
A decision-support system that combines ontology-based knowledge 
representation, semantic similarity between problem descriptions 
and scientific literature, and structured implementation support 
can provide relevant and justified machine learning method 
recommendations for users with limited machine learning expertise.
\end{quote}


To investigate this hypothesis, the following research questions 
are defined:

\begin{itemize}
    \item How can ontology-based knowledge improve machine learning method selection?
    \item How can semantic similarity between problem descriptions and literature support recommendations?
    \item How can literature-based evidence support and justify 
machine learning method recommendations?
    \item To what extent can such a system support early-stage machine learning implementation?
\end{itemize}

To investigate the hypothesis and research questions, a web-based decision-support
system, called MLguide, was developed that combines structured knowledge
with information from scientific literature.

\section{Thesis Structure}
\label{sec:thesis-structure}

The remainder of this thesis is structured as follows:

\begin{itemize}
    \item Chapter 1 introduces the motivation, problem description, and research questions
    \item Chapter 2 presents background and related work
    \item Chapter 3 presents the methodology and system overview
    \item Chapter 4 presents the system implementation
    \item Chapter 5 presents the results
    \item Chapter 6 discusses the findings and limitations
    \item Chapter 7 concludes the thesis and suggests future work
\end{itemize}





